\section{Implications}

If the picture holds, a handful of consequences follow that are not decorative additions but the actual payoff of the construction. Each one ties together facts that are usually treated as separate, and each one falls out of the same discrete growth that defines the substrate. The implications below are stated as readings of the model rather than as independent claims, and where the finite and infinite versions of the cosmos give different stories the text gives both without choosing between them. The choice is left open on purpose, and the model is built so that nothing of importance hangs on it.

\subsection{Time, the arrow, and expansion are one process}

The mesh grows by one beat per shell. That single growth carries three jobs that physics normally keeps apart. It is the time step, since each beat advances the state. It is the arrow of time, since growth has a direction that cannot be reversed without unwinding the whole structure. And it is the spatial expansion, since the new shells are literally more space. Cosmic expansion is not a term added to a static stage. The expansion \emph{is} the universe growing, and time is the growing. Three usually separate facts collapse into one.

\subsection{An age and a horizon, on either reading}

The model does not prefer a finite cosmos over an infinite one. Both readings are live, and this is treated as a genuine open question rather than a settled point. The growth picture pays off either way. On the finite reading the mesh has a true first shell, the seed, and a finite size that acts as a horizon. Both come at no extra cost, since they are just the start and the present edge of the growth. On the infinite or eternal reading there is no first shell. The growth is beginningless and endless. Yet any observer still sees a finite past horizon and a finite size, because only finitely many shells lie inside their light cone. On that reading the horizon is observer relative rather than a birth. Either way one gets a finite observable age and a finite horizon for free. What differs is only whether the very first shell exists, and the model leaves that to the reader.

\subsection{Grand unification as geometry}

The radial direction of the bulk is the renormalization scale, in the manner of holographic renormalization. Read that way, unification becomes geometric rather than imposed. The forces are one structure, an $SO(10)$ symmetry, deep in the bulk. They split into the Standard Model on the flat cusp layer near the boundary. The running of the couplings is then the radial coarse-graining itself, the same motion outward through shells. The grand unified theory is not a separate high-energy postulate. It is the geometry seen at the unification scale, with the splitting into known forces being what that geometry looks like once it has been coarse-grained out to the flat layer.

\subsection{Computability and a discrete base}

The base is discrete and local. Every finite region of it is therefore computable by a finite procedure, with no appeal to a continuum and no uncomputable real numbers hiding at the bottom. The universe is simulable in principle, region by region, to any finite extent. This holds whether the cosmos is finite and growing or infinite and eternal. The load-bearing commitment is not finiteness of size but discreteness of structure. A continuum would smuggle in real numbers of unbounded precision at every point, and the model declines to do that. The cosmos is finite-or-eternal in extent but always discrete in grain.

\subsection{Consciousness as the substrate}

Because the base is experiential, which is the monist commitment of the model, the deepest reading is that the universe is at bottom a structure of consciousness, and that matter, space, gravity, and the forces are shapes that structure takes. This is a reading of the same formal substrate, not an extra ingredient. The tone, the rule, and the growth are unchanged. What changes is only what one takes the substrate to be made of, and the model takes it to be made of experience.

\subsection{Why these implications matter}

The five readings are the return on the construction. Expansion and time become one process. A beginning and a horizon come for free on either the finite or the infinite story. Unification turns into geometry. The cosmos is computable. And the base is experiential. None of these is added by hand. Each is forced by the same discrete growing substrate that the earlier parts build up. The predictions in the next chapter make the package falsifiable, and the open problems chapter says plainly what remains unfinished, so that the implications here are read as consequences with a stated cost rather than as a finished account.
